\documentclass[10pt]{article}
\usepackage[margin=20mm]{geometry}
\usepackage{amsmath,amssymb,mathtools,fontspec,unicode-math,xurl,hyperref}
\IfFontExistsTF{DejaVu Serif}{\setmainfont{DejaVu Serif}}{\setmainfont{Latin Modern Roman}}
\setmathfont{latinmodern-math.otf}
\setlength{\parindent}{0pt}\setlength{\parskip}{6pt}\sloppy
\title{Scientific review of
Correlation decay for finite lattice gauge theories at weak coupling}
\author{Initial scientific review}
\date{}
\begin{document}
\maketitle
Paper. Arka Adhikari and Sky Cao, Annals of Probability 53 (2025), 140–174. DOI: 10.1214/24-AOP1702 (https://doi.org/10.1214/24-AOP1702). The reviewed source is the acquired 35-page journal PDF deposited with NSF. Page references below are printed journal pages; physical PDF page \(p\) corresponds to journal page \(139+p\).\par
\section*{Summary}
The inspected proof reduces correlation decay to a swapping involution and then bounds the event on which the swap is unavailable by a Peierls sum. The preservation of the union of supports makes the involution well defined even though the knot decomposition and its coordinate bijection are chosen rather than unique. The covariance cancellation, the weight bound, and the final numerical threshold survive direct checks.\par
One local proof defect is established. Lemma 6.8 is stated for arbitrary plaquette sets, but its first inference requires a common coordinate separating the entire sets. Pairwise \(\ell^\infty\) separation does not imply that condition. The actual application uses rectangles, for which the required implication is valid. The assessment is therefore Minor, confined to the unsupported generality of that auxiliary proof; no failure of Theorem 1.1 has been established.\par
\section*{Scope, hypotheses, and dependencies}
The setting is a finite four-dimensional cubic lattice and a finite group \(G\), equipped with a unitary representation of character \(\chi\). The energy per nonidentity holonomy is controlled by \[\Delta_G=\min_{g\ne1}\operatorname{Re}\{\chi(1)-\chi(g)\}.\] A nonvacuous finite-temperature threshold requires \(\Delta_G>0\); a faithful unitary representation has this property. A representation with a nontrivial kernel gives \(\Delta_G=0\) and does not satisfy a finite version of the stated threshold. The observables are bounded functions of loops contained in rectangles \(B_i\), invariant under conjugation of each coordinate separately. Their support size \(|B_i|\) counts plaquettes, not vertices or perimeter. Theorem 1.1 assumes \[\beta\Delta_G\ge114+4\log|G|.\] This is a weak-coupling assertion for this finite-group model, without a claim about continuous gauge groups or all temperatures.\par
The review follows Sections 2–6, with particular attention to the constant-fiber passage to homomorphisms, the hypotheses of the gluing maps, the swap, and the counting estimate. Appendix A.1 was checked against the gauge-fixing notation. The argument still relies on the external support-count and topological gluing facts cited from Cao’s work, and on the combinatorial counting input used in Appendix A.3. Those external proofs have not been independently certified here. The Higgs-model extensions and their later appendices are outside the detailed checked scope.\par
\section*{The swapping and Peierls arguments}
Write \(\nu\) for the law on based homomorphisms. For independent \(\Psi_1,\Psi_2\sim\nu\), Lemma 3.2 (p. 148) gives \[|\operatorname{Cov}(h_1(\Psi_1),h_2(\Psi_1))|
 \le 2\|h_1\|_\infty\|h_2\|_\infty\,\mathbb P(E^c).\] Indeed, on \(E\), a bijection preserving the product mass equates the sum of \(h_1(\Psi_1)h_2(\Psi_1)\) with the sum of \(h_1(\Psi_1)h_2(\Psi_2)\). The two remaining sums on \(E^c\) each have the displayed norm bound. This calculation does not presume independence after conditioning on \(E\).\par
In Sections 4–5 (pp. 150–154), gluing is applied to sets separated by a good rectangle; arbitrary disjoint plaquette sets do not suffice. For \[P=\operatorname{supp}\psi_1\cup\operatorname{supp}\psi_2\cup B_1\cup B_2,\] the chosen knot-coordinate map is fixed as a function of \(P\). Swapping the coordinate containing \(B_2\) preserves the union of the two supports. Consequently the second application uses the same \(P\), knot decomposition, and coordinate map, so \(T^2=1\). Product weights are unchanged because the two corresponding component factors are merely interchanged. Independent conjugacy invariance of the loop functions is exactly what is needed to pass from the component holonomies to the original observables.\par
Equation reference: eq:sum. Put \(A=4\cdot10^{24}|G|^2\), \(b=\beta\Delta_G\), and \(m_0=|B_1|+|B_2|\). Lemmas 6.6–6.7 (pp. 156–157) remove a specified knot from the positive partition sum and bound its weight by \(4^m|G|^{2m}e^{-b(m-m_0)}\). This comparison does not assert probabilistic independence of different knots. Combining it with the anchored knot count gives, for separated rectangles and \(L\ge1\), \[
 \mathbb P(E^c)\le A^{m_0}\frac{(Ae^{-b})^{L-1}}{1-Ae^{-b}}.\] For nonintegral \(L\), the integer lower limit is rounded upward, which can only decrease this upper bound. The numerical assumption gives \[Ae^{-b/2}\le4\cdot10^{24}e^{-57}<1,\qquad
 Ae^{-b}\le e^{-b/2}<\tfrac12 .\] Thus [eq:sum] (\#eq:sum) is at most \(2A^{m_0}e^{-b(L-1)/2}\), and the covariance factor supplies the constant \(4\) in Theorem 1.1. The strict numerical inequality was also verified using a rational Taylor lower bound for \(e^{57}\). For \(0\le L\le1\), the generic bound \(2\|h_1\|_\infty\|h_2\|_\infty\) already implies the theorem; the negative-power comparison in [eq:sum] (\#eq:sum) is unnecessary. A rectangle with no plaquettes has at most one nondegenerate coordinate, so all its closed edge paths have trivial holonomy and its observable is constant.\par
\section*{A local overstatement in the geometric lemma}
Problem 1 (The first inference of Lemma 6.8).  Lemma 6.8 (p. 157) assumes that every vertex of a plaquette set \(P_1\) is at \(\ell^\infty\) distance at least \(L\) from every vertex of \(P_2\). Its proof immediately infers a single coordinate direction in which the two entire sets are separated by at least \(L\). This inference is false for general plaquette sets.\par
Verification and limits. Take \(L=3\) and a sufficiently large cubic lattice. Let \(Q_t\) be the unit plaquette whose vertices are \[\{(t,y,z,0): y,z\in\{0,1\}\}.\] Set \(P_1=\{Q_{-3},Q_3\}\) and \(P_2=\{Q_0\}\). Every vertex pair from the two sets has first-coordinate difference \(3\) and all other coordinate differences at most \(1\); hence every such distance is \(3\). Nevertheless, in the first coordinate \(P_1\) lies on both sides of \(P_2\), and in each other coordinate the projections coincide. There is no direction \(i\), and no ordering of the two sets, giving the globally separating inequalities used in the proof. An exact enumeration of these vertices confirms distance \(3\) and zero interval gap in all four coordinates.\par
This example refutes the inferred geometric premise. It does not itself refute the conclusion about the size of a containing knot: that conclusion could require another argument for general sets. The slicing proof as written is valid when a coordinate globally separates the sets. In particular, if \(B_1,B_2\) are axis-aligned rectangles, then \[\min_{x\in B_1,y\in B_2}\|x-y\|_\infty
 =\max_i \operatorname{dist}(I_{1,i},I_{2,i}),\] where \(I_{r,i}\) is the coordinate interval of \(B_r\). Minimization can be performed in every coordinate simultaneously. Therefore separation by \(L\) supplies the required coordinate for the actual application. The slicing hyperplanes then force the \(L-1\) extra plaquettes, with the integer interpretation described above. ◻\par
A sufficient correction for the main proof is to state the needed geometric bound for coordinate-separated sets, or just for the rectangles used in Proposition 6.1, and to treat \(L\le1\) by the trivial covariance estimate. This restores the inspected use, without purporting to prove the more general formulation of Lemma 6.8. The defect has no demonstrated central impact and is assessed as Minor.\par
\section*{Notation clarification and overall assessment}
Definition 2.13 (p. 147) prints the condition \(\sigma_e=1\) for all \(e\in\Lambda_1\) in the definition of \(\mathrm{GF}(T)\). The intended range is \(e\in T\). Literally fixing every edge leaves only the identity configuration and cannot parametrize all homomorphisms. Appendix A.1 (pp. 158–159) explicitly fixes tree edges and assigns the free generators to non-tree edges, establishing the intended bijection. This is a notation correction, not a separately graded mathematical defect.\par
The principal correlation estimate survives the checked algebra and the rectangle-specific geometric argument, conditional on the external topological and counting inputs. The present result is a bounded initial review, not a certification of all extensions or an independent acceptance decision. A bounded search of the authors’ publication information and arXiv record identified the accepted preprint version and the cited topology paper; it did not locate a separate correction or supplement. This does not establish that none exists.\par
Assessment | Repair verification | Author adoption | Remaining scope\par
Minor: generality of Lemma 6.8’s proof | Rectangle case and final constants checked | Unknown | General-set conclusion; external topology; Higgs extensions\par
\section*{Scope and references}
The exact finite diagnostic uses integer vertex coordinates and rational arithmetic; no author code was executed. The raw source parsing, cited-page images, exact-source compilation, dependency inventories, and subsequent report-page observations are recorded separately with their actual file hashes. Scientific argument and page inspection were performed by the assigned initial-review AI session; a native helper’s role was restricted to clerical metadata and file bindings. No human assessment is represented.\par
9 A. Adhikari and S. Cao, Correlation decay for finite lattice gauge theories at weak coupling, Ann. Probab. 53 (2025), 140–174, DOI: 10.1214/24-AOP1702. Reviewed journal PDF: https://par.nsf.gov/servlets/purl/10650878. A. Adhikari and S. Cao, arXiv:2202.10375, public version record, v3 dated 17 March 2024. Consulted for source history, not substituted for the acquired journal text. S. Cao, Wilson loop expectations in lattice gauge theories with finite gauge groups, Comm. Math. Phys. 380 (2020), 1439–1505, DOI: 10.1007/s00220-020-03912-z; arXiv:2001.05627. Bibliographic record and its role as a dependency were inspected; the complete external proof was not independently verified.\par
\end{document}
